% Source PDF: source/普通高考/2016/2016全国1文(河南,河北,山西,江西,湖北,湖南,广东,安徽,福建).pdf, page 1.
% PDF embedded bitmap attempt: pdfimages -list found no embedded bitmap objects; the flowchart is vector line art.
% Grid evidence: tmp/redraw_flowchart_2016_national_paper_1_liberal/grid/page1_grid100.png,
% with q10 positioned from the ungridded page render.
% Elements: rounded start/end terminals, input/output parallelograms, computation and increment boxes,
% decision diamond, directed connectors, loop-back branch, and 是/否 labels.
% Relations: 开始 -> 输入(x,y,n) -> (x,y update) -> x^2+y^2>=36; 是 -> 输出(x,y) -> 结束;
% 否 -> n=n+1 -> back to the computation step.
% solid segments: all box outlines and connector lines; dashed segments: none.
% labels: branch labels are placed beside the corresponding diamond exits; all math text is centered
% with zero paragraph indentation and compact node padding.
\begin{tikzpicture}[
  >=Stealth,
  line width=0.45pt,
  line join=round,
  line cap=round,
  font=\normalsize,
  flow/.style={draw, align=center,
               execute at begin node={\setlength{\parindent}{0pt}},
               inner xsep=5pt, inner ysep=3pt},
  terminal/.style={flow, rounded corners=0.18cm, minimum width=1.18cm, minimum height=0.58cm},
  io/.style={flow, trapezium, trapezium left angle=72, trapezium right angle=108,
             minimum width=2.45cm, minimum height=0.70cm},
  process/.style={flow, minimum width=4.15cm, minimum height=0.78cm},
  increment/.style={flow, minimum width=1.78cm, minimum height=0.70cm},
  decision/.style={flow, diamond, aspect=2.7, minimum width=3.75cm, minimum height=1.05cm,
                  inner xsep=1pt, inner ysep=1pt},
  branch label/.style={execute at begin node={\setlength{\parindent}{0pt}}},
]
  \node[terminal] (start) at (0,7.15) {开始};
  \node[io] (input) at (0,6.02) {输入 $x,y,n$};
  \node[process] (update) at (0,4.75) {$x=x+\dfrac{n-1}{2},y=ny$};
  \node[decision] (test) at (0,3.35) {$x^2+y^2\geqslant36$};
  \node[io] (output) at (0,2.03) {输出 $x,y$};
  \node[terminal] (finish) at (0,0.92) {结束};
  \node[increment] (inc) at (-3.20,4.75) {$n=n+1$};

  \draw[->] (start.south) -- (input.north);
  \draw[->] (input.south) -- (update.north);
  \draw[->] (update.south) -- (test.north);
  \draw[->] (test.south) -- node[branch label, right=2pt] {是} (output.north);
  \draw[->] (output.south) -- (finish.north);

  % 否 branch: decision -> increment, then return to the computation step.
  \draw[->] (test.west) -- (-3.20,3.35) -- (inc.south);
  % Keep the return line centered in the narrow gap between input and update;
  % the final short drop makes the incoming edge explicit without touching input.
  \coordinate (loop-arrow-tip) at (-0.70,5.405);
  \draw[->] (inc.north) -- (-3.20,5.405) -- (loop-arrow-tip);
  \draw (loop-arrow-tip) -- (0,5.405) -- (update.north);
  \node[branch label, anchor=east, inner sep=1pt] at (-1.96,3.15) {否};
\end{tikzpicture}
